% ====================================================================
% FILE: content/10_klevels_full_ru.tex
% Полная иерархия уровней K (K0–K12)
% Онтология континуумов — Core 1.3.3
% ====================================================================

\section{Полная иерархия континуумов \texorpdfstring{\texorpdfstring{$K_0$}{K\_0}–$K_{12}$}{K0–K12}}
\label{sec:klevels-full}

В этой главе представлено полное восстановленное описание всех уровней
континуума \(K_0\)–\(K_{12}\), объединяющее материал, первоначально
распределённый между archival source corpus и трассами Physics, Chemistry, Biology, Cognition
и Social.
В отличие от Раздела~the formal model chapter, где был дан сжатый обзор,
в настоящей главе каждый уровень раскрывается через полное структурное
определение.

Core~1.3.3 фиксирует иерархию как \(K_0\)–\(K_{12}\), поскольку открытый
источник уже содержит явные верхнеуровневые модули, кросс-уровневые переходы
и оболочки предсказания/процесса для \(K_{11}\) и \(K_{12}\).
Если рассматривать иерархию как бы прерывающуюся на \(K_{10}\),
монография становилась бы внутренне противоречивой.

Каждый уровень континуума задаётся кортежем
\[
    K_x = \big(\Omega_x, A_x, P_x(t), J_x(t), \Theta_x,
           \partial\Omega_x, C_x, k_x(t)\big).
\]
Уровни различаются природой осей, потенциалов, порогов,
геометрией границ, циклами и структурой пространства вложения
\(M_x\), необходимого для их существования.


% ====================================================================
\subsection{Обзор вертикальной иерархии}
% ====================================================================

Вертикальная иерархия континуумов монотонна:
\[
    K_0 \rightarrow K_1 \rightarrow K_2 \rightarrow \dots
    \rightarrow K_{10} \rightarrow K_{11} \rightarrow K_{12},
\]
каждый переход соответствует возникновению новых осей,
порогов, потоков и циклов, которые не сводимы к уровням ниже.

Общий обзор показан в Таблице~source label table klevels summary.

\begin{table}[h!]
\centering
\begingroup
\footnotesize
\RaggedRight
\setlength{\tabcolsep}{4pt}
\begin{tabular}{@{}L{0.10\textwidth}L{0.25\textwidth}L{0.49\textwidth}@{}}
\hline
Уровень & Доминирующие оси & Характерная структура \\
\hline
\(K_0\) & Ось различий & Структурная подложка \\
\(K_1\) & Одномерная геометрическая ось & Классический континуум \\
\(K_2\) & Физические полевые оси & Фазы, BKT, масса, когерентность \\
\(K_3\) & Химические оси & Сети RAF, каталитика \\
\(K_4\) & Мембранные оси & Осмотические, кривизны, пороги проницаемости \\
\(K_5\) & Оси возбудимости & Прото-спайки, ионные потоки \\
\(K_6\) & Когнитивные оси & Связывание, внутренние модели \\
\(K_7\) & Нормативные/социальные оси & Институты, доверие \\
\(K_8\) & Цивилизационные оси & Инфраструктура, коллективные циклы \\
\(K_9\) & Теоретические оси & Парадигмы, онтологии \\
\(K_{10}\) & Мета-теоретические оси & Рекурсия моделей моделей \\
\(K_{11}\) & Мета-эволюционные формальные оси & Операторы, мета-логики, развивающиеся пространства моделей \\
\(K_{12}\) & Оси глобальной согласованности & Межфреймовая семантическая и структурная согласованность \\
\hline
\end{tabular}
\endgroup
\caption{Глобальный обзор континуумов \texorpdfstring{\(K_0\)}{K0}–\(K_{12}\).}
\label{tab:klevels-summary}
\end{table}

Каждый уровень требует пространства вложения
\[
   M_0 \subset M_1 \subset \dots \subset M_{10} \subset M_{11} \subset M_{12},
\]
с осями \(A(M_x)\), достаточными для размещения континуумов данного уровня.


% ====================================================================
\subsection{Уровень \texorpdfstring{\texorpdfstring{$K_0$}{K\_0}}{K0}: Структурная подложка}
% ====================================================================

\paragraph{Формальные данные.}
\[
    K_0 = (S,\Delta,\mathcal{C}),
\]

содержит:
\begin{itemize}
    \item \(S\) — множество различимых состояний;
    \item \(\Delta\) — структурная функция различимости;
    \item \(\mathcal{C}\) — структурное отношение (или семейство отношений),
          сохраняющее различимость.
\end{itemize}

На этом уровне отсутствуют время, геометрия, потоки и потенциалы.

\paragraph{Порог.}
Порог существования:
\[
   \Theta_0 = \varepsilon > 0,
   \qquad
   s_1 \neq s_2 \Rightarrow \Delta(s_1,s_2) \ge \varepsilon.
\]

\paragraph{Роль.}
Обеспечивает логическую подложку различимости;
без него не может существовать ни один континуум.


% ====================================================================
\subsection{Уровень \texorpdfstring{\texorpdfstring{$K_1$}{K\_1}}{K1}: Одномерные континуумы}
% ====================================================================

\paragraph{Формальные данные.}
\[
    K_1 = (\Omega_1, A_1, P_1, J_1, \Theta_1,
           \partial\Omega_1, C_1, k_1).
\]

\begin{itemize}
    \item \(A_1\): одна геометрическая ось (1D континуум);
    \item \(\Omega_1\): классическое конфигурационное пространство
          (например, \(C^0(T,H^1(X))\cap C^1(T,L^2(X))\));
    \item \(P_1\): классический энергетический функционал;
    \item \(J_1\): классический динамический поток;
    \item \(\Theta_1\): пороги устойчивости и критических значений.
\end{itemize}

\paragraph{Примеры.}
Классические поля в одной пространственной размерности, связанные осцилляторы,
модели скалярной диффузии.

\paragraph{Циклы.}
Периодические орбиты, устойчивые предельные циклы.


% ====================================================================
\subsection{Уровень \texorpdfstring{\texorpdfstring{$K_2$}{K\_2}}{K2}: Физические континуумы}
% ====================================================================

\paragraph{Оси.}
Многомерные полевые оси:
\begin{itemize}
    \item пространственные оси;
    \item внутренние полевые оси (например, калибровочные степени свободы);
    \item оси параметров порядка (когерентность, намагниченность и т.д.).
\end{itemize}

\paragraph{Потенциалы и потоки.}
Физические энергетические функционалы, уравнения поля, потоки ренормгруппы.

\paragraph{Пороги.}
\begin{itemize}
    \item \(\Theta_{\rm crit}\): границы фаз;
    \item \(\Theta_{\rm BKT}\): порог расплетания вихрей;
    \item \(\Theta_{\rm mass}\): порог когерентности для возникновения массы;
    \item \(\Theta_{\rm death}\): конфайнмент или декуплинг.
\end{itemize}

\paragraph{Циклы.}
Топологические солитоны, пары вихрей, циклы поля.


% ====================================================================
\subsection{Уровни \texorpdfstring{\texorpdfstring{$K_3$}{K\_3}–\texorpdfstring{$K_4$}{K\_4}}{K3–K4}: Химические и протоклеточные континуумы}
% ====================================================================

\subsubsection{Уровень \texorpdfstring{\(K_3\)}{K3}: Химические континуумы}

\paragraph{Оси.}
Оси химических концентраций, оси окружающей среды (pH, солёность, температура).

\paragraph{Пороги.}
\begin{itemize}
    \item порог существования RAF (замыкание);
    \item пороги каталитической активации;
    \item порог перколяции для реакционных сетей.
\end{itemize}

\paragraph{Циклы.}
Химические циклы, циклы замыкания RAF.

\subsubsection{Уровень \texorpdfstring{\(K_4\)}{K4}: Протоклеточные континуумы}

\paragraph{Оси.}
Мембранные оси: кривизна, поверхностное натяжение, проницаемость, заряд.

\paragraph{Пороги.}
\begin{itemize}
    \item осмотический порог \(\Theta_{\rm grad}\);
    \item порог проницаемости \(\Theta_{\rm perm}\);
    \item порог разрыва мембраны \(\Theta_{\rm mem}\);
    \item порог кривизны \(\Theta_{\rm curv}\).
\end{itemize}

\paragraph{Циклы.}
Метаболические циклы, циклы роста мембраны, циклы поддержания градиента.

\paragraph{Геометрия границы.}
Детализированная участковая модель:
состояния patch \(\sigma_i\in\{L_\alpha,L_\beta,L_o,L_f,L_b\}\),
локальные векторы порогов \(\Theta_i\),
взаимодействия на уровне участков и динамика разрыва.


% ====================================================================
\subsection{Уровень \texorpdfstring{\texorpdfstring{$K_5$}{K\_5}}{K5}: Ранние нейронные и биоэлектрические континуумы}
% ====================================================================

\paragraph{Оси.}
Оси возбудимости:
\begin{itemize}
    \item ось потенциала мембраны \(\Delta V\);
    \item оси концентраций ионов;
    \item оси состояний каналов.
\end{itemize}

\paragraph{Потенциалы и потоки.}
\begin{itemize}
    \item \(P_5\): электрохимические потенциалы;
    \item \(J_5\): ионные потоки, потоки переменных ионных каналов.
\end{itemize}

\paragraph{Пороги.}
\begin{itemize}
    \item порог возбуждения \(\Theta_{\rm exc}\);
    \item порог рефрактерности \(\Theta_{\rm refr}\);
    \item градиентный порог \(\Theta_{\rm grad}\);
    \item порог отказа \(\Theta_{\rm death}\) (потеря возбудимости).
\end{itemize}

\paragraph{Циклы.}
Цикл прото-спайков \(C_{\rm spike}\),
ранние осцилляционные циклы,
петли возбуждения уровня patch.


% ====================================================================
\subsection{Уровень \texorpdfstring{\texorpdfstring{$K_6$}{K\_6}}{K6}: Когнитивные континуумы}
% ====================================================================

\paragraph{Оси.}
\begin{itemize}
    \item репрезентативные оси;
    \item оси связывания;
    \item концептуальные оси;
    \item предиктивные оси.
\end{itemize}

\paragraph{Потенциалы.}
\begin{itemize}
    \item потенциалы семантической энергии;
    \item потенциалы ошибки предсказания;
    \item потенциалы стабилизации памяти.
\end{itemize}

\paragraph{Пороги.}
\begin{itemize}
    \item порог связывания \(\Theta_{\rm bind}\);
    \item порог предсказания \(\Theta_{\rm pred}\);
    \item порог когерентности \(\Theta_{\rm coh}\);
    \item порог экспрессии \(\Theta_{\rm expr}\).
\end{itemize}

\paragraph{Циклы.}
Когнитивные циклы:
\begin{itemize}
    \item цикл внимания;
    \item цикл предсказания;
    \item цикл консолидации памяти;
    \item цикл обновления модели.
\end{itemize}


% ====================================================================
\subsection{Уровни \texorpdfstring{\texorpdfstring{$K_7$}{K\_7}–\texorpdfstring{$K_8$}{K\_8}}{K7–K8}: Социальные и цивилизационные континуумы}
% ====================================================================

\subsubsection{Уровень \texorpdfstring{\(K_7\)}{K7}: Социальные континуумы}

\paragraph{Оси.}
\begin{itemize}
    \item нормативные оси;
    \item оси доверия;
    \item рольовые оси;
    \item институциональные оси.
\end{itemize}

\paragraph{Пороги.}
\begin{itemize}
    \item порог доверия \(\Theta_{\rm trust}\);
    \item порог легитимности \(\Theta_{\rm leg}\);
    \item порог рольовой когерентности \(\Theta_{\rm role}\).
\end{itemize}

\paragraph{Циклы.}
Нормативные циклы, институциональные циклы, координационные циклы.

\subsubsection{Уровень \texorpdfstring{\(K_8\)}{K8}: Цивилизационные континуумы}

\paragraph{Оси.}
Масштабные оси:
\begin{itemize}
    \item инфраструктурные оси;
    \item оси потоков энергии;
    \item оси коммуникации;
    \item оси сложности.
\end{itemize}

\paragraph{Пороги.}
\begin{itemize}
    \item порог системного стресса;
    \item порог ёмкости;
    \item порог коезии.
\end{itemize}

\paragraph{Циклы.}
Энергетические циклы, экономические циклы, мега-институциональные циклы.


% ====================================================================
\subsection{Уровни \texorpdfstring{\texorpdfstring{$K_9$}{K\_9}–$K_{12}$}{K9–K12}: Теоретические, мета-теоретические и континуумы верхней согласованности}
% ====================================================================

\subsubsection{Уровень \texorpdfstring{\(K_9\)}{K9}: Теоретические континуумы}

\paragraph{Оси.}
\begin{itemize}
    \item концептуальные оси;
    \item формальные оси;
    \item оси моделирования.
\end{itemize}

\paragraph{Пороги.}
\begin{itemize}
    \item порог непротиворечивости;
    \item порог когерентности;
    \item порог экспрессивности \(\Theta_{\rm expr}\).
\end{itemize}

\paragraph{Циклы.}
Циклы эволюции теории, парадигмальные циклы.

\subsubsection{Уровень \(K_{10}\): Мета-теоретические континуумы}

\paragraph{Оси.}
\begin{itemize}
    \item метаязыковые оси;
    \item функториальные оси;
    \item оси модели моделей.
\end{itemize}

\paragraph{Пороги.}
\begin{itemize}
    \item порог метакогерентности;
    \item порог самоссылки;
    \item порог мета-экспрессивности.
\end{itemize}

\paragraph{Циклы.}
Функториальные циклы, циклы мета-теоретической рекурсии.

\subsubsection{Уровень \(K_{11}\): Рекурсивные мета-теоретические континуумы}

\paragraph{Оси.}
\begin{itemize}
    \item Оси семейств операторов;
    \item оси металогику;
    \item оси преобразования пространств моделей;
    \item оси рекурсивного структурного отражения.
\end{itemize}

\paragraph{Пороги.}
\begin{itemize}
    \item порог когерентности операторов;
    \item порог рекурсивного замыкания;
    \item порог совместимости мета-пространства.
\end{itemize}

\paragraph{Роль.}
\(K_{11}\) необходима, когда теории, наборы операторов и пространства моделей
рассматриваются как эволюционирующие структурные объекты, а не как статические
контейнеры. Поэтому это первый уровень, на котором рамка явно изучает
мета-эволюцию.

\paragraph{Циклы.}
Циклы пересмотра операторов, циклы обновления металогики,
циклы рекурсии пространств моделей.

\subsubsection{Уровень \(K_{12}\): Континуумы глобальной согласованности}

\paragraph{Оси.}
\begin{itemize}
    \item оси семантической согласованности;
    \item оси глобальной совместимости;
    \item оси межфреймового связывания.
\end{itemize}

\paragraph{Пороги.}
\begin{itemize}
    \item порог глобальной непротиворечивости;
    \item порог семантической согласованности;
    \item порог верхней допустимости.
\end{itemize}

\paragraph{Роль.}
\(K_{12}\) — минимально определяемый верхний слой, необходимый для обсуждения
ограничений глобальной согласованности между семействами структур \(K_{11}\).
Его функция структурная и ограничающая: он фиксирует, что у иерархии существует
явная проблема глобальной согласованности на верхнем уровне, а не просто дальнейший
рост размерности нижних уровней.

\paragraph{Циклы.}
Циклы восстановления глобальной согласованности, циклы семантического выравнивания,
циклы межфреймового замыкания.


% ====================================================================
\subsection{Кросс-уровневое резюме}
% ====================================================================

Континуум \(K_x\) существует как действующий континуум тогда и только тогда,
когда
\[
   \Omega_x \neq \emptyset,
   \qquad
   A_x \subseteq A(M_x),
   \qquad
   \\Theta_x \\text{ выполним на } M_x.
\]

Условия вертикальной непрерывности:
\begin{itemize}
    \item \(A_x \subset A_{x+1}\);
    \item \(M_x \subset M_{x+1}\);
    \item \(\Theta_{x+1}\) уточняют \(\Theta_x\);
    \item потоки \(J_x\) вложены в \(J_{x+1}\);
    \item циклы \(C_x\) образуют подструктуры \(C_{x+1}\).
\end{itemize}

Завершение этой главы даёт восстановленную полную иерархию континуумов от
\(K_0\) до \(K_{12}\).
% GENERATED_BY: oc_core_monograph_translation_draft_pipeline_v1.py
% AI_ASSISTED_TRANSLATION_DRAFT: TRUE
% THEOREM_REVIEW_STATUS: PENDING
% THEOREM_REVIEW_MODE: FORMAL_AI_GATE__CODEX_CLI_CHATGPT
% THEOREM_REVIEW_REVIEWER_ID: 
% THEOREM_REVIEW_AUTHORITY_CLASS: 
% THEOREM_REVIEWED_AT: 
% THEOREM_REVIEW_DOSSIER_REF: 
% SOURCE_LANGUAGE: EN
% TARGET_LANGUAGE: RU
% SOURCE_EN_REF: content/10_klevels_full.tex
% ====================================================================
